%=========================
% FrontMatter/summary_vi.tex
% Tóm tắt báo cáo bằng tiếng Việt
%=========================
\begin{center}
    \textbf{\Large TÓM TẮT}
\end{center}

Trong bối cảnh thương mại điện tử ngày càng phát triển mạnh mẽ, nhu cầu tìm kiếm và tư vấn sản phẩm trực tuyến bằng ngôn ngữ tự nhiên hoặc hình ảnh ngày càng trở nên phổ biến. Tuy nhiên, các hệ thống tìm kiếm truyền thống chủ yếu dựa trên từ khóa thường gặp nhiều hạn chế trong việc hiểu đúng nhu cầu người dùng và cung cấp kết quả phù hợp. Vì vậy, việc ứng dụng trí tuệ nhân tạo để xây dựng hệ thống hỗ trợ mua sắm thông minh là một hướng nghiên cứu có tiềm năng và giá trị thực tiễn cao.

Đồ án tập trung nghiên cứu và xây dựng hệ thống \textit{Trợ lý Mua sắm Đa phương tiện sử dụng RAG và Computer Vision} nhằm hỗ trợ người dùng tìm kiếm, tư vấn và gợi ý sản phẩm hiệu quả hơn trong lĩnh vực thương mại điện tử. Hệ thống cho phép người dùng tương tác thông qua hội thoại tự nhiên hoặc tải lên hình ảnh để tìm kiếm các sản phẩm tương tự. Điểm nổi bật của đề tài là việc kết hợp các mô hình ngôn ngữ lớn LLM, kỹ thuật RAG, Vector Database và các mô hình thị giác máy tính nhằm xây dựng hệ thống tìm kiếm đa phương thức thông minh.

Về khía cạnh xử lý ngôn ngữ tự nhiên, đồ án triển khai chatbot tư vấn mua sắm có khả năng hiểu yêu cầu của người dùng, truy xuất dữ liệu sản phẩm và sinh phản hồi phù hợp theo ngữ cảnh hội thoại. Đồng thời, hệ thống áp dụng các mô hình Computer Vision và Multimodal Embedding để trích xuất đặc trưng hình ảnh, hỗ trợ tìm kiếm sản phẩm dựa trên ảnh đầu vào. Toàn bộ hệ thống được tích hợp trên nền tảng Web Fullstack, kết hợp triển khai bằng Docker nhằm đảm bảo khả năng mở rộng, tính linh hoạt và tốc độ phản hồi ổn định.

Cuối cùng, đồ án tiến hành đánh giá hiệu năng hệ thống thông qua các tiêu chí như độ chính xác tìm kiếm, khả năng truy xuất thông tin và thời gian phản hồi của các mô hình AI. Kết quả nghiên cứu góp phần định hướng phát triển các hệ thống trợ lý mua sắm thông minh ứng dụng trí tuệ nhân tạo, giúp nâng cao trải nghiệm người dùng và tối ưu hóa hoạt động tìm kiếm sản phẩm trong lĩnh vực Thương mại điện tử.